% ============================================================
% APÊNDICE A: ESPECIFICAÇÃO DOS CASOS DE USO
% ============================================================
\section*{APÊNDICE A: ESPECIFICAÇÃO DOS CASOS DE USO}
\phantomsection
\addcontentsline{toc}{section}{APÊNDICE A: ESPECIFICAÇÃO DOS CASOS DE USO}

Esta seção apresenta a especificação textual de dez casos de uso que compõem o escopo funcional central da plataforma WeFIND, descrevendo atores, pré-condições, fluxos principais, alternativas e exceções.

\vspace{0.5cm}

\subsection*{Caso de Uso UC01: Efetuar Login}
\noindent
\begin{table}[H]
\centering
\small
\begin{tabular}{|>{\raggedright\arraybackslash}p{3.2cm}|>{\raggedright\arraybackslash}p{11.8cm}|}
\hline
Identificador & UC01 \\ \hline
Nome do Caso de Uso & Efetuar Login \\ \hline
Ator Principal & Usuário Visitante \\ \hline
Atores Secundários & Serviço de Autenticação (Supabase Auth) e Gateway de Mensageria (Evolution API / WhatsApp) \\ \hline
Resumo & Permite ao usuário autenticar-se na plataforma mediante e-mail e senha ou por meio da conta Google, gerando token JWT de sessão persistente. \\ \hline
Pré-Condições & Possuir cadastro ativo no WeFIND. \\ \hline
\multicolumn{2}{|l|}{Fluxo Principal:} \\
\multicolumn{2}{|>{\raggedright\arraybackslash}p{15.4cm}|}{
1. O usuário aciona a opção ``Entrar'' na tela inicial.\par
2. O sistema exibe o formulário solicitando e-mail e senha.\par
3. O usuário preenche as credenciais e toca em ``Entrar''.\par
4. O sistema valida o formato sintático do e-mail e o comprimento da senha.\par
5. O sistema submete a requisição ao Supabase Auth.\par
6. O serviço valida as credenciais e retorna a sessão autenticada.\par
7. O sistema armazena a sessão localmente (\textit{AsyncStorage}) e redireciona ao mural principal autenticado.
} \\ \hline
\multicolumn{2}{|l|}{Fluxos Alternativos:} \\
\multicolumn{2}{|>{\raggedright\arraybackslash}p{15.4cm}|}{
3a. Esquecimento de credenciais (\textit{<<extend>>}): o usuário aciona a opção Redefinir Senha, recebendo código OTP de 6 dígitos via Evolution API no WhatsApp para cadastro de nova senha.\par
3b. Login com Google: o usuário aciona a opção de entrar com o Google, conclui a autenticação no navegador do dispositivo e retorna ao aplicativo com a sessão gerenciada pelo Supabase Auth.\par
1a. Exploração sem cadastro: o usuário opta por ``Continuar explorando como visitante'', acessando a visualização de publicações e do mapa em modo somente leitura.
} \\ \hline
\multicolumn{2}{|l|}{Fluxos de Exceção:} \\
\multicolumn{2}{|>{\raggedright\arraybackslash}p{15.4cm}|}{
4a. Campos obrigatórios vazios: o sistema impede o envio e exibe avisos de preenchimento obrigatório nos campos destacados.\par
6a. Credenciais incorretas: o sistema emite alerta visual de erro e retém o usuário no formulário para nova tentativa.
} \\ \hline
Pós-Condições & Sessão iniciada com permissões de publicação e comunicação ativas. \\ \hline
\end{tabular}
\end{table}

\vspace{0.8cm}

\subsection*{Caso de Uso UC02: Realizar Cadastro}
\noindent
\begin{table}[H]
\centering
\small
\begin{tabular}{|>{\raggedright\arraybackslash}p{3.2cm}|>{\raggedright\arraybackslash}p{11.8cm}|}
\hline
Identificador & UC02 \\ \hline
Nome do Caso de Uso & Realizar Cadastro \\ \hline
Ator Principal & Usuário Visitante \\ \hline
Atores Secundários & Supabase Auth e Gateway de Mensageria (Evolution API / WhatsApp) \\ \hline
Resumo & Permite a novos membros criarem uma conta informando dados cadastrais e realizando a verificação por código de segurança no WhatsApp. \\ \hline
Pré-Condições & Dispositivo conectado à internet. \\ \hline
\multicolumn{2}{|l|}{Fluxo Principal:} \\
\multicolumn{2}{|>{\raggedright\arraybackslash}p{15.4cm}|}{
1. O usuário toca na opção ``Cadastre-se'' na tela de autenticação.\par
2. O sistema abre o formulário de registro solicitando nome completo, e-mail, telefone WhatsApp e senha.\par
3. O usuário preenche as informações e confirma o cadastro.\par
4. O sistema valida os campos, cria o registro no Supabase Auth e aciona por inclusão (\textit{«include»}) o caso de uso Verificar Código, disparando um código OTP de 6 dígitos via Evolution API para o WhatsApp informado.\par
5. O sistema exibe a tela de confirmação de telefone com 6 campos numéricos.\par
6. O usuário digita o código recebido no WhatsApp.\par
7. O caso de uso Verificar Código valida a sequência numérica e ativa a conta de usuário permanentemente.
} \\ \hline
\multicolumn{2}{|l|}{Fluxos Alternativos:} \\
\multicolumn{2}{|>{\raggedright\arraybackslash}p{15.4cm}|}{
6a. Reenvio de código de segurança: caso não receba o código, o usuário aciona o botão de reenvio após o término da contagem regressiva do temporizador de segurança.
} \\ \hline
\multicolumn{2}{|l|}{Fluxos de Exceção:} \\
\multicolumn{2}{|>{\raggedright\arraybackslash}p{15.4cm}|}{
4a. E-mail ou telefone já cadastrados: o sistema emite alerta indicando a duplicidade cadastral e impede o registro repetido.\par
4b. Senha fora dos parâmetros: o sistema sinaliza que a senha deve possuir no mínimo 6 caracteres.\par
6b. Código incorreto ou expirado: o sistema sinaliza a inconsistência e solicita nova inserção.
} \\ \hline
Pós-Condições & Conta criada, com telefone confirmado por código e pronta para autenticação. \\ \hline
\end{tabular}
\end{table}

\clearpage

\subsection*{Caso de Uso UC03: Manter Publicação}
\noindent
\begin{table}[H]
\centering
\small
\begin{tabular}{|>{\raggedright\arraybackslash}p{3.2cm}|>{\raggedright\arraybackslash}p{11.8cm}|}
\hline
Identificador & UC03 \\ \hline
Nome do Caso de Uso & Manter Publicação \\ \hline
Ator Principal & Usuário (Autenticado) \\ \hline
Atores Secundários & Banco de Dados (PostgreSQL) e Supabase Storage \\ \hline
Resumo & Permite cadastrar, editar, encerrar como reencontrado ou excluir anúncios de animais perdidos e encontrados. \\ \hline
Pré-Condições & Usuário autenticado e com permissão de localização ativa. \\ \hline
\multicolumn{2}{|l|}{Fluxo Principal:} \\
\multicolumn{2}{|>{\raggedright\arraybackslash}p{15.4cm}|}{
1. O usuário aciona o botão de nova publicação (``+'') na barra de navegação.\par
2. O sistema apresenta o formulário solicitando a categoria: ``Perdido'' ou ``Encontrado''.\par
3. O usuário seleciona a categoria correspondente à situação do animal.\par
4. O usuário adiciona fotografias (com recorte integrado) e seleciona os atributos físicos do pet (espécie, raça, porte, cor, idade).\par
5. O sistema executa obrigatoriamente por inclusão (\textit{<<include>>}) o caso de uso Selecionar Localização no Mapa, capturando as coordenadas geográficas.\par
6. O usuário confirma a publicação.\par
7. O sistema envia as fotos ao Storage e grava os dados na tabela \textit{items}.\par
8. A publicação passa a ser exibida no mural e no mapa para a comunidade.
} \\ \hline
\multicolumn{2}{|l|}{Fluxos Alternativos:} \\
\multicolumn{2}{|>{\raggedright\arraybackslash}p{15.4cm}|}{
4a. Edição de publicação existente: o autor atualiza fotos ou descrições do anúncio já publicado.\par
4b. Encerramento com sucesso: o autor altera o status para ``Reencontrado'', encerrando a busca.
} \\ \hline
\multicolumn{2}{|l|}{Fluxos de Exceção:} \\
\multicolumn{2}{|>{\raggedright\arraybackslash}p{15.4cm}|}{
4c. Imagem inválida ou tamanho excessivo: o sistema impede o envio e solicita a seleção de uma foto válida.\par
5a. Ponto no mapa não demarcado: o sistema bloqueia a submissão e exige a marcação das coordenadas geográficas.
} \\ \hline
Pós-Condições & Publicação persistida, visível no mapa e disponível para a comunidade. \\ \hline
\end{tabular}
\end{table}

\vspace{0.8cm}

\subsection*{Caso de Uso UC04: Registrar Avistamento}
\noindent
\begin{table}[H]
\centering
\small
\begin{tabular}{|>{\raggedright\arraybackslash}p{3.2cm}|>{\raggedright\arraybackslash}p{11.8cm}|}
\hline
Identificador & UC04 \\ \hline
Nome do Caso de Uso & Registrar Avistamento \\ \hline
Ator Principal & Usuário (Autenticado) \\ \hline
Atores Secundários & Banco de Dados (PostgreSQL) e Serviço de Notificações \\ \hline
Resumo & Permite a um membro da comunidade informar novas coordenadas e fotos do local onde viu um animal perdido na rua, atualizando o rastro de busca. \\ \hline
Pré-Condições & O animal deve estar publicado com status ``Perdido''. \\ \hline
\multicolumn{2}{|l|}{Fluxo Principal:} \\
\multicolumn{2}{|>{\raggedright\arraybackslash}p{15.4cm}|}{
1. Ao visualizar os detalhes de um animal perdido, o usuário aciona a opção ``Vi esse Pet''.\par
2. O sistema abre o modal de avistamento solicitando foto do local e relato.\par
3. O usuário captura ou anexa uma fotografia do animal avistado.\par
4. O sistema inclui obrigatoriamente (\textit{<<include>>}) o caso de uso Selecionar Localização no Mapa, capturando o ponto GPS do avistamento.\par
5. O usuário confirma o envio do avistamento.\par
6. O sistema atualiza o marcador principal do animal para as novas coordenadas e grava o histórico em \textit{sightings}.\par
7. O sistema registra uma notificação para o autor da publicação informando o novo avistamento e, quando houver permissão do dispositivo, envia um aviso ao celular.
} \\ \hline
\multicolumn{2}{|l|}{Fluxos Alternativos:} \\
\multicolumn{2}{|>{\raggedright\arraybackslash}p{15.4cm}|}{
4a. Ajuste manual de localização: caso o sinal de GPS esteja impreciso, o usuário arrasta manualmente o pino até o local visualizado na rua.
} \\ \hline
\multicolumn{2}{|l|}{Fluxos de Exceção:} \\
\multicolumn{2}{|>{\raggedright\arraybackslash}p{15.4cm}|}{
4b. Coordenadas geográficas ausentes: o sistema impede a gravação do avistamento enquanto o ponto no mapa não for confirmado.\par
5a. Falha temporária de rede: o sistema retém os dados digitados e sugere nova tentativa de envio.
} \\ \hline
Pós-Condições & Marcador do animal reposicionado no mapa, tutor alertado e histórico de rastro atualizado. \\ \hline
\end{tabular}
\end{table}

\clearpage

\subsection*{Caso de Uso UC05: Enviar Mensagem}
\noindent
\begin{table}[H]
\centering
\small
\begin{tabular}{|>{\raggedright\arraybackslash}p{3.2cm}|>{\raggedright\arraybackslash}p{11.8cm}|}
\hline
Identificador & UC05 \\ \hline
Nome do Caso de Uso & Enviar Mensagem \\ \hline
Ator Principal & Usuário (Autenticado) \\ \hline
Atores Secundários & Serviço de Mensagens em Tempo Real (Supabase Realtime) \\ \hline
Resumo & Viabiliza canal direto e privativo de troca de mensagens e imagens entre o tutor do animal e voluntários da comunidade sem expor contatos telefônicos. \\ \hline
Pré-Condições & Ambos os usuários devem possuir cadastro ativo. \\ \hline
\multicolumn{2}{|l|}{Fluxo Principal:} \\
\multicolumn{2}{|>{\raggedright\arraybackslash}p{15.4cm}|}{
1. Na tela de detalhes da publicação, o usuário aciona o botão ``Entrar em Contato''.\par
2. O sistema cria ou recupera a sala de conversa existente entre os dois usuários para aquele animal.\par
3. O sistema abre a interface de chat em tempo real.\par
4. O usuário digita a mensagem ou anexa uma foto e confirma o envio.\par
5. O sistema transmite a mensagem em tempo real e notifica o destinatário pela central de notificações e, quando houver permissão, no celular.\par
6. A mensagem é persistida na tabela \textit{messages} e exibida na conversa.
} \\ \hline
\multicolumn{2}{|l|}{Fluxos Alternativos:} \\
\multicolumn{2}{|>{\raggedright\arraybackslash}p{15.4cm}|}{
5a. Destinatário desconectado: a mensagem é gravada de modo permanente no banco de dados e sincronizada assim que o destinatário abrir o aplicativo.
} \\ \hline
\multicolumn{2}{|l|}{Fluxos de Exceção:} \\
\multicolumn{2}{|>{\raggedright\arraybackslash}p{15.4cm}|}{
1a. Tentativa de contato consigo mesmo: o sistema oculta a opção de chat quando o visualizador é o próprio autor do anúncio.\par
4a. Mensagem vazia: o sistema desabilita o botão de envio enquanto não houver conteúdo textual ou anexo de mídia.
} \\ \hline
Pós-Condições & Canal de comunicação estabelecido com histórico mantido na caixa de entrada. \\ \hline
\end{tabular}
\end{table}

\vspace{0.8cm}

\subsection*{Caso de Uso UC06: Compartilhar Publicação}
\noindent
\begin{table}[H]
\centering
\small
\begin{tabular}{|>{\raggedright\arraybackslash}p{3.2cm}|>{\raggedright\arraybackslash}p{11.8cm}|}
\hline
Identificador & UC06 \\ \hline
Nome do Caso de Uso & Compartilhar Publicação \\ \hline
Ator Principal & Usuário / Usuário Visitante \\ \hline
Atores Secundários & Módulo Gerador de Código QR e API Nativa de Compartilhamento \\ \hline
Resumo & Permite a qualquer usuário renderizar automaticamente artes gráficas diagramadas em múltiplos formatos (A4, formatos verticais e murais de redes sociais) contendo código QR para divulgação rápida. \\ \hline
Pré-Condições & O animal deve possuir ao menos uma fotografia cadastrada. \\ \hline
\multicolumn{2}{|l|}{Fluxo Principal:} \\
\multicolumn{2}{|>{\raggedright\arraybackslash}p{15.4cm}|}{
1. Na tela de detalhes da publicação, o usuário clica em ``Compartilhar Cartaz''.\par
2. O sistema abre o modal com pré-visualização ao vivo da arte gerada.\par
3. No formato A4, o sistema gera o código QR que direciona à página pública da publicação.\par
4. O usuário seleciona o formato desejado (A4 para impressão, proporção vertical 9:16 ou proporção quadrada 1:1).\par
5. O usuário escolhe entre salvar a imagem na galeria ou acionar o compartilhamento direto via WhatsApp ou redes sociais.\par
6. O sistema exporta o cartaz como arquivo de imagem.
} \\ \hline
\multicolumn{2}{|l|}{Fluxos Alternativos:} \\
\multicolumn{2}{|>{\raggedright\arraybackslash}p{15.4cm}|}{
4a. Cópia de texto rápido: o sistema oferece botão auxiliar para copiar texto formatado contendo os dados essenciais e o link web do animal para colar em conversas.
} \\ \hline
\multicolumn{2}{|l|}{Fluxos de Exceção:} \\
\multicolumn{2}{|>{\raggedright\arraybackslash}p{15.4cm}|}{
1a. Ausência de fotografias: o sistema informa que é obrigatório dispor de pelo menos uma foto para diagramar o cartaz de busca.
} \\ \hline
Pós-Condições & Cartaz digital gerado para impressão física ou compartilhamento em redes sociais. \\ \hline
\end{tabular}
\end{table}

\clearpage

\subsection*{Caso de Uso UC07: Manter Moderação e Denúncias}
\noindent
\begin{table}[H]
\centering
\small
\begin{tabular}{|>{\raggedright\arraybackslash}p{3.2cm}|>{\raggedright\arraybackslash}p{11.8cm}|}
\hline
Identificador & UC07 \\ \hline
Nome do Caso de Uso & Manter Moderação e Denúncias \\ \hline
Ator Principal & Administrador \\ \hline
Atores Secundários & Banco de Dados (PostgreSQL) \\ \hline
Resumo & Permite ao administrador inspecionar relatórios de denúncia enviados pela comunidade sobre anúncios indevidos, aplicando sanções ou exclusão do conteúdo. \\ \hline
Pré-Condições & Conta autenticada com perfil de administrador (\textit{isAdmin = true}). \\ \hline
\multicolumn{2}{|l|}{Fluxo Principal:} \\
\multicolumn{2}{|>{\raggedright\arraybackslash}p{15.4cm}|}{
1. O administrador acessa a tela de moderação do sistema.\par
2. O sistema lista os anúncios que receberam denúncias dos usuários com o respectivo motivo.\par
3. O administrador examina as fotos, o texto e os dados do anúncio reportado.\par
4. O administrador seleciona a ação cabível: desconsiderar denúncia, suspender anúncio ou excluir publicação.\par
5. O sistema atualiza o status do anúncio e da denúncia conforme a decisão tomada.
} \\ \hline
\multicolumn{2}{|l|}{Fluxos Alternativos:} \\
\multicolumn{2}{|>{\raggedright\arraybackslash}p{15.4cm}|}{
4a. Descarte de denúncia infundada: o administrador constata a conformidade do anúncio e arquiva o chamado mantendo a publicação ativa.\par
4b. Solicitação de ajuste: o administrador suspende temporariamente o anúncio e solicita alterações de dados ao usuário.
} \\ \hline
\multicolumn{2}{|l|}{Fluxos de Exceção:} \\
\multicolumn{2}{|>{\raggedright\arraybackslash}p{15.4cm}|}{
4c. Infração grave ou fraude dolosa: o administrador exclui o anúncio e suspende preventivamente a conta do usuário infrator.
} \\ \hline
Pós-Condições & Denúncia analisada e publicação mantida, suspensa ou removida conforme a decisão do administrador. \\ \hline
\end{tabular}
\end{table}

\clearpage

\clearpage

\subsection*{Caso de Uso UC08: Solicitar Comprovação de Tutela}
\noindent
\begin{table}[H]
\centering
\small
\begin{tabular}{|>{\raggedright\arraybackslash}p{3.2cm}|>{\raggedright\arraybackslash}p{11.8cm}|}
\hline
Identificador & UC08 \\ \hline
Nome do Caso de Uso & Solicitar Comprovação de Tutela \\ \hline
Ator Principal & Usuário Autenticado (solicitante) \\ \hline
Atores Secundários & Autor da publicação (quem acolheu o animal) \\ \hline
Resumo & Permite ao usuário que se identifica como tutor de um animal anunciado como encontrado comprovar a tutela para obter o endereço exato e combinar a devolução. \\ \hline
Pré-Condições & Usuário autenticado; publicação do tipo ``Encontrado'' ativa, com endereço exato protegido. \\ \hline
\multicolumn{2}{|l|}{Fluxo Principal:} \\
\multicolumn{2}{|>{\raggedright\arraybackslash}p{15.4cm}|}{
1. Na ficha do animal encontrado, o usuário aciona a opção de comprovação de tutela.\par
2. O sistema solicita a forma de comprovação (fotos do animal, carteira de vacinação ou marcas ocultas).\par
3. O usuário anexa as imagens pela galeria ou pela câmera e escreve uma mensagem ao autor da publicação.\par
4. O usuário confirma o envio.\par
5. O sistema registra a solicitação como pendente e notifica o autor da publicação por notificação e mensagem privada.\par
6. O autor analisa o relato e as imagens e aprova a solicitação.\par
7. O sistema libera ao solicitante o endereço exato do local onde o animal está acolhido e o notifica da aprovação.
} \\ \hline
\multicolumn{2}{|l|}{Fluxos Alternativos:} \\
\multicolumn{2}{|>{\raggedright\arraybackslash}p{15.4cm}|}{
6a. Recusa da solicitação: o autor recusa o pedido e pode registrar o motivo; o sistema notifica o solicitante, e o endereço permanece protegido.
} \\ \hline
\multicolumn{2}{|l|}{Fluxos de Exceção:} \\
\multicolumn{2}{|>{\raggedright\arraybackslash}p{15.4cm}|}{
3a. Mensagem vazia: o sistema exige a descrição de como o usuário comprova ser o tutor antes de permitir o envio.
} \\ \hline
Pós-Condições & Solicitação registrada e analisada; em caso de aprovação, endereço exato liberado ao solicitante. \\ \hline
\end{tabular}
\end{table}


\clearpage
\subsection*{Caso de Uso UC09: Manter Meus Pets}
\noindent
\begin{table}[H]
\centering
\small
\begin{tabular}{|>{\raggedright\arraybackslash}p{3.2cm}|>{\raggedright\arraybackslash}p{11.8cm}|}
\hline
Identificador & UC09 \\ \hline
Nome do Caso de Uso & Manter Meus Pets \\ \hline
Ator Principal & Usuário Autenticado (tutor) \\ \hline
Atores Secundários & Usuários próximos, que recebem o alerta quando o animal é declarado perdido \\ \hline
Resumo & Permite ao tutor cadastrar seus animais de estimação de forma preventiva, gerar a carteirinha digital de cada um e, se o animal desaparecer, transformar o cadastro em uma publicação de animal perdido. \\ \hline
Pré-Condições & Usuário autenticado. \\ \hline
\multicolumn{2}{|l|}{Fluxo Principal:} \\
\multicolumn{2}{|>{\raggedright\arraybackslash}p{15.4cm}|}{
1. No perfil, o usuário abre a área ``Meus pets e carteirinhas''.\par
2. O usuário toca em adicionar e informa foto, nome, espécie, raça, porte, cor, sexo, idade ou data de nascimento, microchip, castração, vacinação e observações médicas.\par
3. O sistema salva o animal na lista do tutor.\par
4. O usuário toca em ``Ver carteirinha''.\par
5. O sistema gera a carteirinha digital com os dados do animal, o nome e o contato do tutor e um código QR, que pode ser salva ou compartilhada.
} \\ \hline
\multicolumn{2}{|l|}{Fluxos Alternativos:} \\
\multicolumn{2}{|>{\raggedright\arraybackslash}p{15.4cm}|}{
4a. Animal desaparecido: o usuário aciona ``Meu pet fugiu'' e confirma; o sistema cria uma publicação de animal perdido com as características já preenchidas, avisa os usuários próximos e compara a publicação com os anúncios de animais encontrados.\par
4b. Edição ou exclusão: o usuário altera os dados do animal ou o remove da lista.
} \\ \hline
\multicolumn{2}{|l|}{Fluxos de Exceção:} \\
\multicolumn{2}{|>{\raggedright\arraybackslash}p{15.4cm}|}{
3a. Falha ao salvar: o sistema informa que não foi possível salvar os dados e mantém o usuário no formulário.
} \\ \hline
Pós-Condições & Animal cadastrado na lista do tutor; no fluxo 4a, publicação de animal perdido ativa no mapa e no mural. \\ \hline
\end{tabular}
\end{table}

\clearpage
\subsection*{Caso de Uso UC10: Disponibilizar para Adoção}
\noindent
\begin{table}[H]
\centering
\small
\begin{tabular}{|>{\raggedright\arraybackslash}p{3.2cm}|>{\raggedright\arraybackslash}p{11.8cm}|}
\hline
Identificador & UC10 \\ \hline
Nome do Caso de Uso & Disponibilizar para Adoção \\ \hline
Ator Principal & Usuário Autenticado (autor da publicação) \\ \hline
Atores Secundários & Usuários interessados na adoção \\ \hline
Resumo & Permite oferecer um animal para adoção responsável, seja em uma publicação criada para esse fim, seja em uma publicação de animal encontrado que já passou pelo período mínimo de busca pelo tutor. \\ \hline
Pré-Condições & Usuário autenticado; para animais encontrados, publicação ativa com o animal sob os cuidados do autor. \\ \hline
\multicolumn{2}{|l|}{Fluxo Principal:} \\
\multicolumn{2}{|>{\raggedright\arraybackslash}p{15.4cm}|}{
1. Na ficha da própria publicação de um animal encontrado, o autor aciona ``Disponibilizar para Adoção''.\par
2. O sistema verifica se a publicação completou o período mínimo de sete dias de busca pelo tutor.\par
3. O sistema pede confirmação ao autor.\par
4. O autor confirma.\par
5. O sistema marca a publicação como disponível para adoção, e ela passa a aparecer como ``Para Adoção'' no mural.\par
6. Os interessados entram em contato com o autor pelo chat da publicação.
} \\ \hline
\multicolumn{2}{|l|}{Fluxos Alternativos:} \\
\multicolumn{2}{|>{\raggedright\arraybackslash}p{15.4cm}|}{
1a. Adoção direta: no cadastro de uma nova publicação, o usuário escolhe a opção ``Para Adoção''; a publicação já é criada como disponível para adoção.\par
5a. Pausar adoção: o autor pode retirar o animal da listagem de adoção a qualquer momento.
} \\ \hline
\multicolumn{2}{|l|}{Fluxos de Exceção:} \\
\multicolumn{2}{|>{\raggedright\arraybackslash}p{15.4cm}|}{
2a. Período de busca não concluído: o sistema informa quantos dias faltam e mantém a publicação ativa como animal encontrado.
} \\ \hline
Pós-Condições & Publicação disponível para adoção no mural, com contato pelo chat. \\ \hline
\end{tabular}
\end{table}

\clearpage

% ============================================================
% APÊNDICE B: INSTRUMENTOS DE PESQUISA
% ============================================================
\section*{APÊNDICE B: INSTRUMENTOS DE PESQUISA (QUESTIONÁRIO DE CAMPO E ESCALA SUS)}
\phantomsection
\addcontentsline{toc}{section}{APÊNDICE B: INSTRUMENTOS DE PESQUISA (QUESTIONÁRIO DE CAMPO E ESCALA SUS)}

Este apêndice documenta o instrumento de coleta de dados primários aplicado durante o diagnóstico exploratório (Google Forms) e o protocolo SUS aplicado na avaliação de usabilidade descrita no Capítulo 8.

\subsection*{B.1 Levantamento de Campo com Tutores (Gráficos Consolidados e Instrumento de Pesquisa)}

Esta seção documenta os resultados consolidados da pesquisa de campo aplicada entre 1º e 11 de setembro de 2026 aos participantes voluntários na plataforma Google Forms, durante a etapa de diagnóstico empírico e levantamento de requisitos (Capítulo~3). A divulgação ocorreu em um grupo do Facebook relacionado à proteção animal. Até 11 de setembro, 20 pessoas haviam respondido ao formulário, composto por 18 perguntas. Quando indicado no texto ou nas figuras, uma questão incluída posteriormente possui 19 respostas válidas. As Figuras~B1 a B14 apresentam as distribuições percentuais das respostas obtidas nas dimensões investigadas.

% Configuração da numeração de figuras do Apêndice B (Figura B1, B2, ...)
\renewcommand{\thefigure}{B\arabic{figure}}
\renewcommand{\theHfigure}{B\arabic{figure}}
\setcounter{figure}{0}

Para mapear o perfil do público participante e compreender a convivência direta com animais de estimação na amostra, o questionário investigou quais espécies domésticas são tuteladas pelos respondentes, conforme ilustrado na Figura~\ref{fig:forms_relacao_animais}.

\begin{figure}[H]
    \centering
    \caption{Relação dos participantes com animais de estimação e espécies tuteladas}
    \label{fig:forms_relacao_animais}
    \includegraphics[width=0.68\textwidth]{imagens/01_pesquisa/forms_relacao_animais.png}
    \fonte{Dados da pesquisa (Google Forms, 2026)}
\end{figure}

Conforme demonstrado na Figura~\ref{fig:forms_relacao_animais} (pergunta de múltipla escolha), 80\% dos respondentes (16 voluntários) são tutores de cães e 70\% (14 voluntários) tutelam gatos, com outros 20\% (4 voluntários) cuidando de espécies adicionais como aves e equinos. Como a questão foi incluída posteriormente e não continha a opção ``não possuo animal'', seus resultados devem ser interpretados apenas como um retrato das respostas válidas obtidas nessa pergunta, sem permitir concluir sobre a totalidade da amostra.

\clearpage

Com o objetivo de identificar os principais gargalos e frustrações vivenciados durante a busca por animais perdidos nos canais convencionais, os participantes foram consultados sobre suas maiores dificuldades, conforme apresentado na Figura~\ref{fig:forms_dificuldades}.

\begin{figure}[H]
    \centering
    \caption{Maiores dificuldades encontradas na divulgação de animais perdidos}
    \label{fig:forms_dificuldades}
    \includegraphics[width=0.68\textwidth]{imagens/01_pesquisa/forms_dificuldades_divulgacao.png}
    \fonte{Dados da pesquisa (Google Forms, 2026)}
\end{figure}

A apuração da Figura~\ref{fig:forms_dificuldades} aponta que 100\% dos participantes (19 respostas) sofrem com postagens que somem rapidamente das linhas do tempo nas redes sociais. Além disso, 89,5\% (17 respostas) apontaram a falta de mapas para saber se o pet foi visto por perto e a precariedade de informações morfológicas nas publicações. O medo de divulgar telefones pessoais foi relatado por 73,7\% dos tutores, e 42,1\% mencionaram a demora na montagem de cartazes e o desconhecimento sobre se a busca continua em andamento.

\clearpage

A busca por localização no mapa é um dos recursos centrais do WeFIND. Para mensurar o valor percebido dessa abordagem, questionou-se aos participantes o nível de utilidade de acompanhar publicações em um mapa interativo com filtro de proximidade, como mostra a Figura~\ref{fig:forms_utilidade_mapa}.

\begin{figure}[H]
    \centering
    \caption{Nível de utilidade de visualizar animais perdidos em mapa interativo}
    \label{fig:forms_utilidade_mapa}
    \includegraphics[width=0.68\textwidth]{imagens/01_pesquisa/forms_utilidade_mapa.png}
    \fonte{Dados da pesquisa (Google Forms, 2026)}
\end{figure}

Os dados da Figura~\ref{fig:forms_utilidade_mapa} indicam avaliação positiva de todos os respondentes: 85\% (17 respostas) atribuíram nota máxima 5 e os restantes 15\% (3 respostas) conferiram nota 4 ao recurso de mapa georreferenciado. Nenhuma avaliação inferior foi assinalada. Esse resultado indica que, para o grupo consultado, a visualização em mapa é percebida como útil para orientar as buscas.

\clearpage

A tempestividade é determinante para o resgate de animais nas primeiras horas após o sumiço. O formulário avaliou o grau de relevância atribuído ao recebimento de alertas imediatos no smartphone segmentados por proximidade, conforme apresentado na Figura~\ref{fig:forms_relevancia_alertas}.

\begin{figure}[H]
    \centering
    \caption{Relevância do recebimento de alertas no celular por bairro ou região}
    \label{fig:forms_relevancia_alertas}
    \includegraphics[width=0.68\textwidth]{imagens/01_pesquisa/forms_relevancia_alertas.png}
    \fonte{Dados da pesquisa (Google Forms, 2026)}
\end{figure}

Conforme indicado na Figura~\ref{fig:forms_relevancia_alertas}, 90\% dos participantes (18 respostas) atribuíram nota máxima (5) à relevância de receber notificações no celular quando um animal desaparecer nas proximidades. Outros 5\% atribuíram nota 4 e 5\% nota 3, sem avaliações negativas (notas 1 ou 2). O resultado foi considerado na definição do requisito de notificações (RF12).

\clearpage

Para verificar se os cidadãos colaborariam ativamente mesmo sem condições de acolher fisicamente um animal desorientado na rua, avaliou-se a intenção de registrar avistamentos rápidos com foto e coordenadas de GPS, conforme ilustrado na Figura~\ref{fig:forms_registro_avistamento}.

\begin{figure}[H]
    \centering
    \caption{Disposição dos participantes para registrar avistamento no mapa comunitário}
    \label{fig:forms_registro_avistamento}
    \includegraphics[width=0.68\textwidth]{imagens/01_pesquisa/forms_registro_avistamento.png}
    \fonte{Dados da pesquisa (Google Forms, 2026)}
\end{figure}

Os resultados expostos na Figura~\ref{fig:forms_registro_avistamento} indicam disposição dos participantes em colaborar: 90\% dos entrevistados (18 voluntários) afirmaram com certeza que utilizariam a opção para registrar avistamentos no mapa e 10\% (2 voluntários) declararam que provavelmente utilizariam. Essa intenção declarada foi considerada no desenvolvimento do registro de avistamentos (RF09).

\clearpage

A confecção manual de cartazes de busca costuma consumir tempo precioso dos tutores em momentos de estresse. O questionário averiguou o impacto da geração automatizada de cartazes padronizados prontos para impressão e redes sociais, como exibido na Figura~\ref{fig:forms_cartaz_automatico}.

\begin{figure}[H]
    \centering
    \caption{Importância da geração automatizada de cartaz padronizado com código QR}
    \label{fig:forms_cartaz_automatico}
    \includegraphics[width=0.68\textwidth]{imagens/01_pesquisa/forms_cartaz_automatico.png}
    \fonte{Dados da pesquisa (Google Forms, 2026)}
\end{figure}

A Figura~\ref{fig:forms_cartaz_automatico} indica a relevância atribuída à geração automática de cartazes: 80\% dos participantes (16 voluntários) avaliaram com nota máxima 5 e 20\% (4 voluntários) com nota 4 a disponibilização de cartazes automáticos com fotos e código QR dinâmico. Não houve avaliações neutras ou negativas, o que reforçou a inclusão desse recurso (RF15).

\clearpage

A segurança na troca de informações entre voluntários e tutores é um aspecto sensível na internet. A pesquisa avaliou a preferência dos usuários por um canal de mensagens instantâneas privativo dentro do aplicativo, conforme demonstrado na Figura~\ref{fig:forms_chat_tempo_real}.

\begin{figure}[H]
    \centering
    \caption{Avaliação da comunicação por chat privativo em tempo real sem expor telefone}
    \label{fig:forms_chat_tempo_real}
    \includegraphics[width=0.68\textwidth]{imagens/01_pesquisa/forms_chat_tempo_real.png}
    \fonte{Dados da pesquisa (Google Forms, 2026)}
\end{figure}

A Figura~\ref{fig:forms_chat_tempo_real} mostra avaliação favorável ao chat interno: 75\% dos respondentes (15 participantes) atribuíram a nota máxima 5 e 20\% (4 participantes) a nota 4 à possibilidade de dialogar com quem encontrou o pet sem expor telefones particulares. Com 95\% de notas 4 ou 5 e apenas 5\% com nota 3, o recurso foi incluído nos requisitos centrais (RF11).

\clearpage

O questionário também investigou a preferência dos participantes sobre a exposição de endereços residenciais em páginas públicas, como detalhado na Figura~\ref{fig:forms_ocultacao_endereco}.

\begin{figure}[H]
    \centering
    \caption{Preferência sobre a ocultação do endereço residencial exato em publicações}
    \label{fig:forms_ocultacao_endereco}
    \includegraphics[width=0.68\textwidth]{imagens/01_pesquisa/forms_ocultacao_endereco.png}
    \fonte{Dados da pesquisa (Google Forms, 2026)}
\end{figure}

Na Figura~\ref{fig:forms_ocultacao_endereco}, 80\% dos entrevistados (16 respostas) manifestaram preferência por manter o endereço residencial e o número predial ocultos do público geral. Outros 10\% declararam não se importar e 10\% mostraram-se indiferentes. Os resultados descrevem as respostas do grupo consultado.

\clearpage

O questionário também consultou os participantes sobre a apresentação de informações que auxiliem a verificar a responsabilidade por um animal, conforme ilustrado na Figura~\ref{fig:forms_comprovacao_tutela}.

\begin{figure}[H]
    \centering
    \caption{Validação da exigência de comprovação de tutela para resgate e prevenção de fraudes}
    \label{fig:forms_comprovacao_tutela}
    \includegraphics[width=0.68\textwidth]{imagens/01_pesquisa/forms_comprovacao_tutela.png}
    \fonte{Dados da pesquisa (Google Forms, 2026)}
\end{figure}

Na Figura~\ref{fig:forms_comprovacao_tutela}, 70\% dos participantes aprovaram a apresentação de evidências documentais ou fotos antigas antes da restituição do animal (45\% destacaram a prevenção a golpes e 25\% apoiaram a medida desde que o fluxo fosse simples). Outros 25\% manifestaram receio de burocratização e 5\% não souberam opinar. Esse resultado fundamentou a inclusão do fluxo de comprovação de tutela (RF18), com envio simples de relato e fotos, para atender também aos participantes preocupados com a burocratização.

\clearpage

O formulário também investigou o interesse dos participantes em temas relacionados à adoção e ao acolhimento temporário, conforme exposto na Figura~\ref{fig:forms_adocao_lar_temporario}.

\begin{figure}[H]
    \centering
    \caption{Interesse na utilização de seção dedicada para adoção e cadastro de lar temporário}
    \label{fig:forms_adocao_lar_temporario}
    \includegraphics[width=0.68\textwidth]{imagens/01_pesquisa/forms_adocao_lar_temporario.png}
    \fonte{Dados da pesquisa (Google Forms, 2026)}
\end{figure}

Na Figura~\ref{fig:forms_adocao_lar_temporario}, 70\% dos voluntários (14 respostas) afirmaram que utilizariam uma seção tanto para adoção quanto para oferecer lar temporário, e 20\% (4 respostas) manifestaram interesse em adoção. Esse resultado orientou a inclusão da adoção responsável no aplicativo (RF20); o cadastro de lar temporário não foi incluído no escopo deste trabalho.

\clearpage

O instrumento de pesquisa também avaliou a opinião dos participantes sobre a possibilidade de oferecer uma recompensa, conforme mostra a Figura~\ref{fig:forms_recompensa_resgate}.

\begin{figure}[H]
    \centering
    \caption{Receptividade à inclusão de recompensa pelo resgate diretamente no anúncio}
    \label{fig:forms_recompensa_resgate}
    \includegraphics[width=0.68\textwidth]{imagens/01_pesquisa/forms_recompensa_resgate.png}
    \fonte{Dados da pesquisa (Google Forms, 2026)}
\end{figure}

Na Figura~\ref{fig:forms_recompensa_resgate}, 95\% dos participantes (19 respostas) apoiaram a possibilidade de oferecer recompensa financeira diretamente na publicação, enquanto 5\% indicaram que poderiam utilizá-la dependendo do caso. Esse resultado fundamentou a inclusão da recompensa opcional nas publicações de animais perdidos (RF04), exibida nos detalhes do anúncio e utilizável como filtro de busca (RF07).

\clearpage

Em cenários de emergência causados pelo sumiço de um pet, cada minuto despendido em cadastros extensos retarda o início das buscas. A pesquisa consultou a expectativa dos participantes quanto ao tempo máximo tolerado para criar e publicar um alerta, como detalhado na Figura~\ref{fig:forms_tempo_publicacao_alerta}.

\begin{figure}[H]
    \centering
    \caption{Expectativa de tempo para criação e publicação de um alerta de busca}
    \label{fig:forms_tempo_publicacao_alerta}
    \includegraphics[width=0.68\textwidth]{imagens/01_pesquisa/forms_tempo_publicacao_alerta.png}
    \fonte{Dados da pesquisa (Google Forms, 2026)}
\end{figure}

Os resultados da Figura~\ref{fig:forms_tempo_publicacao_alerta} indicam que 85\% dos respondentes (17 participantes) esperam concluir o cadastro de um alerta em menos de 2 minutos, e os demais 15\% aceitam entre 2 e 5 minutos para o preenchimento de características completas. Nenhum respondente considerou o tempo irrelevante, consolidando o requisito de formulários com preenchimento assistido e etapas simplificadas.

\clearpage

Para validar o potencial de adesão prática da solução antes de sua disponibilização nas lojas oficiais, o questionário avaliou a disposição voluntária dos entrevistados em instalar e utilizar a aplicação em seus smartphones, conforme apresentado na Figura~\ref{fig:forms_intencao_instalacao}.

\begin{figure}[H]
    \centering
    \caption{Disposição dos participantes para instalar o aplicativo no smartphone}
    \label{fig:forms_intencao_instalacao}
    \includegraphics[width=0.68\textwidth]{imagens/01_pesquisa/forms_intencao_instalacao.png}
    \fonte{Dados da pesquisa (Google Forms, 2026)}
\end{figure}

A Figura~\ref{fig:forms_intencao_instalacao} atesta a ampla receptividade da comunidade: 100\% dos participantes que responderam à questão declararam de maneira categórica que instalariam o aplicativo no aparelho, seja para apoiar resgates na vizinhança ou para proteger seus próprios pets. A ausência de respostas contrárias ou indecisas referenda a viabilidade social da plataforma WeFIND.

\clearpage

Por fim, a pesquisa buscou identificar quais recursos são vistos como os maiores atrativos do WeFIND em comparação com os meios tradicionais de divulgação (como redes sociais e aplicativos de mensagens), conforme exposto na Figura~\ref{fig:forms_diferencial_plataforma}.

\begin{figure}[H]
    \centering
    \caption{Funcionalidade apontada como maior diferencial em relação às redes sociais}
    \label{fig:forms_diferencial_plataforma}
    \includegraphics[width=0.68\textwidth]{imagens/01_pesquisa/forms_diferencial_plataforma.png}
    \fonte{Dados da pesquisa (Google Forms, 2026)}
\end{figure}

A apuração da Figura~\ref{fig:forms_diferencial_plataforma} aponta a centralização temática exclusiva como o principal diferencial para 45\% dos votantes (9 respostas), seguida pelo mapa interativo por proximidade (20\%, 4 respostas), pela emissão automática de cartazes padronizados (15\%, 3 respostas), pelo registro de avistamentos por GPS (10\%, 2 respostas) e pela privacidade do chat interno (10\%, 2 respostas). Esses dados expressam a percepção do grupo consultado e foram considerados na priorização das funcionalidades.



Para documentar o instrumento aplicado, as Figuras~\ref{fig:questionario_parte1} a~\ref{fig:questionario_parte5} reproduzem as perguntas do formulário de pesquisa de campo, na forma como foram apresentadas aos participantes no Google Forms.

\begin{figure}[H]
    \centering
    \caption{Questionário da pesquisa de campo no Google Forms -- parte 1 de 5}
    \label{fig:questionario_parte1}
    \includegraphics[width=0.85\textwidth,height=18cm,keepaspectratio]{imagens/01_pesquisa/questionario_forms_parte1.png}
    \fonte{Autoria própria (Google Forms, 2026)}
\end{figure}

A primeira parte reúne as perguntas sobre identificação dos animais e experiências anteriores de perda. A segunda parte, apresentada na Figura~\ref{fig:questionario_parte2}, trata dos canais de divulgação utilizados, das dificuldades encontradas e da avaliação inicial do mapa, dos alertas e do registro de avistamentos.

\begin{figure}[H]
    \centering
    \caption{Questionário da pesquisa de campo no Google Forms -- parte 2 de 5}
    \label{fig:questionario_parte2}
    \includegraphics[width=0.85\textwidth,height=18cm,keepaspectratio]{imagens/01_pesquisa/questionario_forms_parte2.png}
    \fonte{Autoria própria (Google Forms, 2026)}
\end{figure}

A Figura~\ref{fig:questionario_parte3} reúne a terceira parte do questionário, com perguntas sobre o cartaz automático, o chat privado, a exibição do endereço, a comprovação de tutela e a seção de adoção e lar temporário.

\begin{figure}[H]
    \centering
    \caption{Questionário da pesquisa de campo no Google Forms -- parte 3 de 5}
    \label{fig:questionario_parte3}
    \includegraphics[width=0.85\textwidth,height=18cm,keepaspectratio]{imagens/01_pesquisa/questionario_forms_parte3.png}
    \fonte{Autoria própria (Google Forms, 2026)}
\end{figure}

A quarta parte do formulário, apresentada na Figura~\ref{fig:questionario_parte4}, aborda a recompensa no anúncio, os filtros considerados essenciais, o tempo esperado para publicar um alerta e a intenção de instalar o aplicativo.

\begin{figure}[H]
    \centering
    \caption{Questionário da pesquisa de campo no Google Forms -- parte 4 de 5}
    \label{fig:questionario_parte4}
    \includegraphics[width=0.85\textwidth,height=18cm,keepaspectratio]{imagens/01_pesquisa/questionario_forms_parte4.png}
    \fonte{Autoria própria (Google Forms, 2026)}
\end{figure}

Por fim, a Figura~\ref{fig:questionario_parte5} mostra a última pergunta do questionário, sobre a funcionalidade considerada o maior diferencial do WeFIND em relação às redes sociais.

\begin{figure}[H]
    \centering
    \caption{Questionário da pesquisa de campo no Google Forms -- parte 5 de 5}
    \label{fig:questionario_parte5}
    \includegraphics[width=0.85\textwidth,height=18cm,keepaspectratio]{imagens/01_pesquisa/questionario_forms_parte5.png}
    \fonte{Autoria própria (Google Forms, 2026)}
\end{figure}

\clearpage

\subsection*{B.2 Escala de Usabilidade do Sistema (SUS)}

Escala de avaliação de usabilidade aplicada aos 12 participantes do teste do WeFIND apresentado no Capítulo 8. O formulário adotou a redação em português das dez afirmações da \textit{System Usability Scale} (SUS) conforme apresentada por Teixeira (2015), reproduzida a seguir. A fonte é uma publicação de divulgação sobre a escala e não é citada como estudo de validação psicométrica da tradução.

O formulário de avaliação foi organizado em quatro seções: (1) confirmação do perfil do respondente como tutor de animal de estimação; (2) pergunta única sobre a conclusão das cinco tarefas do roteiro, com as alternativas ``concluí todas as tarefas com facilidade'', ``concluí todas, mas tive dúvida em alguma etapa'', ``concluí apenas algumas tarefas'' e ``não consegui concluir as tarefas''; (3) as dez afirmações da SUS reproduzidas a seguir; e (4) três perguntas abertas, de resposta opcional, sobre os pontos positivos, as dificuldades encontradas e as sugestões de melhoria.

\vspace{0.3cm}
\noindent
Instruções ao participante: Para cada uma das afirmações a seguir, assinale com um ``X'' o nível de concordância que melhor representa sua experiência prática com o aplicativo WeFIND, em uma escala de 1 (Discordo Totalmente) a 5 (Concordo Totalmente):

\vspace{0.5cm}
\begin{table}[H]
\centering
\small
\begin{tabular}{|c|>{\raggedright\arraybackslash}p{8.8cm}|c|c|c|c|c|}
\hline
Nº & Afirmação Avaliada & 1 & 2 & 3 & 4 & 5 \\ \hline
1 & Eu acho que gostaria de usar esse sistema com frequência. & [ ] & [ ] & [ ] & [ ] & [ ] \\ \hline
2 & Eu acho o sistema desnecessariamente complexo. & [ ] & [ ] & [ ] & [ ] & [ ] \\ \hline
3 & Eu achei o sistema fácil de usar. & [ ] & [ ] & [ ] & [ ] & [ ] \\ \hline
4 & Eu acho que precisaria de ajuda de uma pessoa com conhecimentos técnicos para usar o sistema. & [ ] & [ ] & [ ] & [ ] & [ ] \\ \hline
5 & Eu acho que as várias funções do sistema estão muito bem integradas. & [ ] & [ ] & [ ] & [ ] & [ ] \\ \hline
6 & Eu acho que o sistema apresenta muita inconsistência. & [ ] & [ ] & [ ] & [ ] & [ ] \\ \hline
7 & Eu imagino que as pessoas aprenderão como usar esse sistema rapidamente. & [ ] & [ ] & [ ] & [ ] & [ ] \\ \hline
8 & Eu achei o sistema atrapalhado de usar. & [ ] & [ ] & [ ] & [ ] & [ ] \\ \hline
9 & Eu me senti confiante ao usar o sistema. & [ ] & [ ] & [ ] & [ ] & [ ] \\ \hline
10 & Eu precisei aprender várias coisas novas antes de conseguir usar o sistema. & [ ] & [ ] & [ ] & [ ] & [ ] \\ \hline
\end{tabular}
\end{table}

\clearpage

% ============================================================
% APÊNDICE C: GUIA DE UTILIZAÇÃO DO SISTEMA WeFIND (MANUAL DO USUÁRIO)
% ============================================================
\section*{APÊNDICE C: GUIA DE UTILIZAÇÃO DO SISTEMA WeFIND (MANUAL DO USUÁRIO)}
\phantomsection
\addcontentsline{toc}{section}{APÊNDICE C: GUIA DE UTILIZAÇÃO DO SISTEMA WeFIND (MANUAL DO USUÁRIO)}

Este guia prático foi elaborado para orientar tutores de animais domésticos, protetores independentes e membros da comunidade na utilização do aplicativo WeFIND. Organizado em formato de manual de instruções passo a passo, o documento apresenta capturas das telas do sistema e explica os principais fluxos de busca e divulgação de animais perdidos e encontrados.

Para utilizar a ferramenta, o usuário necessita de um smartphone com sistema operacional Android ou iOS e conexão com a internet, seja por rede móvel celular ou Wi-Fi. Durante o uso, o aplicativo pode solicitar autorização para acessar a localização geográfica e a câmera ou galeria de fotos. Esses recursos apoiam a consulta de publicações próximas e o registro de fotografias.

O roteiro está estruturado em torno de cinco procedimentos da aplicação: a criação e gestão da conta; a consulta do mapa e das rotas; o cadastro de animais perdidos ou encontrados; a geração e o compartilhamento de cartazes; e a comunicação privada entre usuários.

Ao longo das páginas seguintes, cada procedimento é ilustrado com imagens das interfaces correspondentes, acompanhadas de orientações sobre os campos a preencher e as ações disponíveis.

Visando resguardar a privacidade e a segurança dos participantes, o sistema não exibe o endereço residencial exato nem o número de telefone particular em páginas públicas, mantendo o diálogo restrito ao canal de mensagens privativo do aplicativo. Caso surjam dúvidas adicionais durante a navegação prática, o usuário pode recorrer à seção de ajuda e suporte presente no menu de configurações da aplicação, que disponibiliza orientações complementares sobre guarda responsável e políticas comunitárias de uso.

\clearpage

% Configuração da numeração de figuras e âncoras do Apêndice C (Figura C1, C2, ...)
\renewcommand{\thefigure}{C\arabic{figure}}
\renewcommand{\theHfigure}{C\arabic{figure}}
\setcounter{figure}{0}

\subsection*{Passo 1: Autenticação e Criação de Conta}
\noindent
Ao iniciar o aplicativo WeFIND pela primeira vez, o cidadão visualiza a tela de autenticação, onde pode se identificar caso já possua conta ativa ou optar por ingressar na comunidade. Para novos participantes, o acesso inicia-se tocando na opção ``Cadastre-se'', que conduz ao formulário unificado de criação de conta. As interfaces correspondentes são apresentadas na Figura~\ref{fig:manual_login_cadastro}.

\begin{figure}[H]
    \centering
    \caption{Telas de Autenticação e Cadastro de Usuário}
    \label{fig:manual_login_cadastro}
    \vspace{0.15cm}
    \begin{minipage}[t]{0.48\textwidth}
        \centering
        \capturaTela{imagens/04_telas_aplicativo/tela-login-v2.jpeg}
        \par\vspace{0.12cm}
        {\small Figura C1.a) Formulário de Autenticação}
    \end{minipage}
    \hfill
    \begin{minipage}[t]{0.48\textwidth}
        \centering
        \capturaTela{imagens/04_telas_aplicativo/tela-criar-conta-v2.jpeg}
        \par\vspace{0.12cm}
        {\small Figura C1.b) Criação de Nova Conta}
    \end{minipage}
    \par\vspace{0.1cm}
    \fonte{Autoria própria (2026)}
\end{figure}

Para acessar a plataforma, o usuário informa seu e-mail e senha cadastrados na tela de autenticação (Figura~C1.a), com a opção de alternar a visibilidade da senha ou solicitar código de recuperação via WhatsApp caso tenha esquecido suas credenciais. Para novos membros, o formulário de cadastro (Figura~C1.b) orienta o preenchimento de nome completo, e-mail, telefone de contato e senha. Ao confirmar o cadastro, o usuário recebe um código de verificação pelo WhatsApp e, após informá-lo, é direcionado ao painel principal de publicações.

\clearpage

\subsection*{Passo 2: Navegação no Mapa Interativo e Traçado de Rotas}
\noindent
Após o ingresso no aplicativo, o usuário pode consultar no mapa as publicações de animais perdidos e encontrados. A Figura~\ref{fig:manual_mapa_rotas} apresenta o mapa principal com marcadores e uma rota até o local associado à publicação selecionada.

\begin{figure}[H]
    \centering
    \caption{Telas de Mapa Interativo e Traçado de Rota até o Animal}
    \label{fig:manual_mapa_rotas}
    \vspace{0.15cm}
    \begin{minipage}[t]{0.48\textwidth}
        \centering
        \capturaTela{imagens/04_telas_aplicativo/tela-mapa-principal.jpeg}
        \par\vspace{0.12cm}
        {\small Figura C2.a) Mapa com Marcadores}
    \end{minipage}
    \hfill
    \begin{minipage}[t]{0.48\textwidth}
        \centering
        \capturaTela{imagens/04_telas_aplicativo/tela-mapa-rota-animal.jpeg}
        \par\vspace{0.12cm}
        {\small Figura C2.b) Traçado de Rota ao Vivo}
    \end{minipage}
    \par\vspace{0.1cm}
    \fonte{Autoria própria (2026)}
\end{figure}

A visualização da Figura~C2 evidencia a convergência entre geolocalização e mobilidade urbana. Na Figura~C2.a, o mapa renderiza marcadores georreferenciados diferenciados por cores e ícones conforme a espécie (cães, gatos, aves e outros animais domésticos), circundados por um raio de busca semitransparente ajustável pelo usuário. Ao tocar sobre qualquer marcador, o sistema exibe um cartão rápido com foto, nome e distância em quilômetros. Ao acionar o botão de navegação, a interface exibida na Figura~C2.b traça automaticamente a rota viária otimizada até as coordenadas do animal, informando o tempo estimado de deslocamento e a distância a percorrer para apoiar a ação imediata do voluntário ou tutor.

\clearpage

\subsection*{Passo 3: Registro de Publicação de Animal Encontrado na Rua}
\noindent
Ao avistar ou resgatar um animal desorientado em vias públicas, o cidadão inicia o fluxo de cadastro comunitário acionando o botão central ``+'' na barra inferior de navegação. O assistente orienta o preenchimento por etapas sequenciais, associando características físicas, fotografias e coordenadas geográficas exatas. A Figura~\ref{fig:manual_cadastro_publicacao} detalha a etapa de caracterização do animal e o georreferenciamento no mapa interativo.

\begin{figure}[H]
    \centering
    \caption{Registro de Publicação e Fixação de Ponto no Mapa}
    \label{fig:manual_cadastro_publicacao}
    \vspace{0.15cm}
    \begin{minipage}[t]{0.48\textwidth}
        \centering
        \capturaTela{imagens/04_telas_aplicativo/tela-cadastrar-encontrado-1.jpeg}
        \par\vspace{0.12cm}
        {\small Figura C3.a) Características Físicas e Fotos}
    \end{minipage}
    \hfill
    \begin{minipage}[t]{0.48\textwidth}
        \centering
        \capturaTela{imagens/04_telas_aplicativo/tela-cadastrar-encontrado-3.jpeg}
        \par\vspace{0.12cm}
        {\small Figura C3.b) Fixação no Mapa}
    \end{minipage}
    \par\vspace{0.1cm}
    \fonte{Autoria própria (2026)}
\end{figure}

O procedimento de registro documentado na Figura~C3 busca organizar as informações disponibilizadas à comunidade. Na Figura~C3.a, o usuário define a categoria da publicação (animal encontrado na rua), anexa fotografias capturadas na hora pela câmera ou selecionadas da galeria, e seleciona espécie, sexo, porte e cores predominantes da pelagem. Na etapa seguinte, exibida na Figura~C3.b, o sistema captura a posição atual do aparelho via GPS ou permite ao colaborador ajustar manualmente o pino sobre o endereço exato onde o animal foi visto, acionando rotinas de geocodificação reversa para preencher automaticamente o endereço antes da publicação no mural geral.

\clearpage

\subsection*{Passo 4: Consulta de Detalhes e Geração de Cartazes com Código QR}
\noindent
Para permitir a identificação precisa de um animal anunciado, o sistema disponibiliza uma tela completa de detalhes com fotografias, características físicas e opções de compartilhamento. A partir dessa tela, o usuário pode gerar cartazes digitais de busca. A Figura~\ref{fig:manual_detalhes_cartaz} apresenta a visualização detalhada da publicação e a peça gráfica de divulgação gerada automaticamente.

\begin{figure}[H]
    \centering
    \caption{Tela de Detalhes do Animal e Cartaz Digital de Busca}
    \label{fig:manual_detalhes_cartaz}
    \vspace{0.15cm}
    \begin{minipage}[t]{0.48\textwidth}
        \centering
        \capturaTela{imagens/04_telas_aplicativo/tela-detalhes-do-animal.jpeg}
        \par\vspace{0.12cm}
        {\small Figura C4.a) Ficha da Publicação}
    \end{minipage}
    \hfill
    \begin{minipage}[t]{0.48\textwidth}
        \centering
        \capturaTela{imagens/04_telas_aplicativo/figura21_tela_cartaz.jpeg}
        \par\vspace{0.12cm}
        {\small Figura C4.b) Cartaz com Código QR}
    \end{minipage}
    \par\vspace{0.1cm}
    \fonte{Autoria própria (2026)}
\end{figure}

A articulação entre a ficha descritiva e a divulgação externa visualiza-se na Figura~C4. Na Figura~C4.a, a tela de detalhes consolida a galeria fotográfica com paginação deslizante, a descrição detalhada de marcas características e cicatrizes, o mapa de localização aproximada e os botões de ação comunitária (``Vi o Pet'', ``Entrar em Contato'' e ``Compartilhar Cartaz''). Ao tocar na opção de compartilhamento, o componente nativo renderiza a tela exibida na Figura~C4.b, diagramando automaticamente o cartaz de busca com fotografia destacada, dados essenciais e, no formato A4, código QR que direciona à página pública da publicação; também há versões nos formatos 9:16 e 1:1 para redes sociais. Como o cartaz pode conter dados informados pelo autor, seu compartilhamento deve considerar essas informações.

\clearpage

\subsection*{Passo 5: Comunicação Privada via Chat Interno}
\noindent
O contato entre quem perdeu e quem encontrou um animal exige um canal direto que não obrigue os participantes a expor seus telefones. Para isso, o WeFIND disponibiliza mensagens privadas em tempo real. A Figura~\ref{fig:manual_mensagens_chat} apresenta a central de mensagens ativas e a conversa privativa em tempo real.

\begin{figure}[H]
    \centering
    \caption{Central de Mensagens e Bate-papo Privativo em Tempo Real}
    \label{fig:manual_mensagens_chat}
    \vspace{0.15cm}
    \begin{minipage}[t]{0.48\textwidth}
        \centering
        \capturaTela{imagens/04_telas_aplicativo/tela-mensagens.jpeg}
        \par\vspace{0.12cm}
        {\small Figura C5.a) Caixa de Mensagens}
    \end{minipage}
    \hfill
    \begin{minipage}[t]{0.48\textwidth}
        \centering
        \capturaTela{imagens/04_telas_aplicativo/tela-conversa.jpeg}
        \par\vspace{0.12cm}
        {\small Figura C5.b) Chat em Tempo Real}
    \end{minipage}
    \par\vspace{0.1cm}
    \fonte{Autoria própria (2026)}
\end{figure}

O sistema de mensageria (Figura~C5) permite o contato privado entre os membros da comunidade. Na central de mensagens (Figura~C5.a), o usuário acompanha todas as conversas em andamento agrupadas por publicação, com identificação visual do pet e contador de mensagens pendentes. Ao selecionar uma conversa (Figura~C5.b), o usuário troca mensagens e fotos com a outra pessoa em tempo real, sem precisar informar o telefone.

\clearpage

\clearpage

% ============================================================
% APÊNDICE D: MATRIZ DE RASTREABILIDADE DE REQUISITOS E VERIFICAÇÕES
% ============================================================
\section*{APÊNDICE D: MATRIZ DE RASTREABILIDADE DE REQUISITOS E VERIFICAÇÕES}
\phantomsection
\addcontentsline{toc}{section}{APÊNDICE D: MATRIZ DE RASTREABILIDADE DE REQUISITOS E VERIFICAÇÕES}

Este apêndice apresenta a rastreabilidade entre as necessidades identificadas no Capítulo~3, os requisitos funcionais centrais do Capítulo~4 e as funcionalidades correspondentes no aplicativo.

A Quadro~D1 também registra a verificação funcional manual relatada pelo autor. As verificações foram realizadas sem roteiro formal previamente documentado; por isso, a matriz não as apresenta como casos de teste padronizados nem como validação independente. O procedimento e suas limitações são descritos no Capítulo~8.

\vspace{0.3cm}
{
\phantomsection\hypertarget{quadro:rastreabilidade}{}\label{quadro:rastreabilidade}
\small
\renewcommand{\arraystretch}{1.18}
\setlength{\tabcolsep}{2.5pt}
\begin{longtable}{|>{\raggedright\arraybackslash}p{4.1cm}|>{\centering\arraybackslash}p{1.8cm}|>{\raggedright\arraybackslash}p{5.4cm}|>{\raggedright\arraybackslash}p{2.7cm}|}
\hline
\multicolumn{4}{|c|}{Quadro D1: Matriz de Rastreabilidade do Sistema WeFIND} \\
\multicolumn{4}{|c|}{(Necessidades $\rightarrow$ Requisitos $\rightarrow$ Funcionalidades $\rightarrow$ Verificação)} \\ \hline
Necessidade identificada (pesquisa ou diagnóstico) & Requisito & Funcionalidade Implementada & Verificação manual \\ \hline
\endfirsthead

\hline
\multicolumn{4}{|c|}{Quadro D1: Matriz de Rastreabilidade do Sistema WeFIND (Continuação)} \\
\multicolumn{4}{|c|}{(Necessidades $\rightarrow$ Requisitos $\rightarrow$ Funcionalidades $\rightarrow$ Verificação)} \\ \hline
Necessidade identificada (pesquisa ou diagnóstico) & Requisito & Funcionalidade Implementada & Verificação manual \\ \hline
\endhead

\hline
\multicolumn{4}{|r|}{\textit{Continua na próxima página}} \\ \hline
\endfoot

\hline
\multicolumn{4}{|c|}{Fonte: Autoria própria (2026)} \\ \hline
\endlastfoot
Pesquisa de campo indicou uso recorrente de redes sociais e necessidade de participação vinculada a uma conta. & RF01 & Cadastro, autenticação, recuperação de senha e encerramento de sessão com Supabase Auth. & Fluxo principal verificado manualmente; sem erro observado. \\ \hline
Pesquisa apontou preocupação com exposição de telefones pessoais; dados básicos de contato precisam ser atualizáveis. & RF02 & Edição dos dados básicos e de contato do perfil do usuário. & Fluxo principal verificado manualmente; sem erro observado. \\ \hline
Os meios atuais foram considerados insuficientes para reunir as informações necessárias sobre o animal. & RF03 & Consulta de publicações em lista e acesso a detalhes, imagens, características e local registrado. & Fluxo principal verificado manualmente; sem erro observado. \\ \hline
Participantes relataram recorrer a redes sociais para divulgar animais perdidos. & RF04 & Cadastro de publicação de animal perdido ou encontrado com descrição, características, imagens, localização e, para animais perdidos, recompensa opcional. & Fluxo principal verificado manualmente; sem erro observado. \\ \hline
Uma publicação precisa poder ser corrigida e encerrada quando a busca termina. & RF05 & Edição, encerramento e exclusão da própria publicação pelo autor. & Fluxo principal verificado manualmente; sem erro observado. \\ \hline
17 dos 19 respondentes indicaram ausência de mapas interativos como dificuldade nos meios atuais. & RF06 & Mapa interativo com marcadores das publicações de animais perdidos e encontrados. & Fluxo principal verificado manualmente; sem erro observado. \\ \hline
A busca em canais generalistas não oferece filtros próprios por características e distância. & RF07 & Pesquisa e filtros por espécie, situação, características, oferta de recompensa e raio de proximidade. & Fluxo principal verificado manualmente; sem erro observado. \\ \hline
A localização associada à publicação deve orientar a procura em campo. & RF08 & Exibição de rota até o local informado na publicação, quando disponível. & Fluxo principal verificado manualmente; sem erro observado. \\ \hline
Relatos comunitários podem fornecer pistas sobre onde e quando o animal foi visto. & RF09 & Registro de avistamento vinculado à publicação, com local, data/horário, descrição e imagem opcional. & Fluxo principal verificado manualmente; sem erro observado. \\ \hline
Relatos incorretos ou desatualizados precisam ser consultáveis e corrigíveis pelo autor. & RF10 & Linha do tempo de avistamentos e edição ou remoção do relato pelo seu autor. & Fluxo principal verificado manualmente; sem erro observado. \\ \hline
14 dos 19 respondentes relataram preocupação com a exposição de telefones pessoais. & RF11 & Conversa privada entre usuários associada à publicação, com envio e recebimento de mensagens. & Fluxo principal verificado manualmente; sem erro observado. \\ \hline
Todos os 19 respondentes apontaram perda de alcance das publicações após algumas horas. & RF12 & Central de notificações para novos avistamentos, mensagens, correspondências, publicações próximas e lembretes relacionados às publicações do usuário. & Fluxo principal verificado manualmente; sem erro observado. \\ \hline
A comunidade precisa poder sinalizar conteúdo suspeito para análise. & RF13 & Envio de denúncia vinculada à publicação ou ao conteúdo reportado. & Fluxo principal verificado manualmente; sem erro observado. \\ \hline
Denúncias exigem análise e ação administrativa para apoiar a integridade das informações. & RF14 & Consulta de denúncias e moderação do conteúdo pelo administrador. & Fluxo principal verificado manualmente; sem erro observado. \\ \hline
80\% dos respondentes avaliaram com nota máxima e 20\% com nota 4 a geração automatizada de cartazes. & RF15 & Geração de cartaz digital com informações da publicação e código QR para exportação ou compartilhamento. & Incluído na tarefa 3 da avaliação com participantes; resultado individual não registrado. \\ \hline
Publicações sem atualização podem permanecer visíveis após o fim da busca. & RF16 & Vigência de 30 dias, aviso três dias antes do vencimento, ocultação após expirar, renovação durante 30 dias de carência e remoção definitiva, após esse prazo, pela rotina local para publicações do usuário autenticado. & Verificação individual não registrada. \\ \hline
A comparação entre anúncios perdidos e encontrados pode apoiar a identificação de possíveis correspondências. & RF17 & Comparação automática de publicações por espécie, características físicas e distância, com apresentação de sugestões de correspondência e confirmação final pelos usuários. & Teste com 20 pares fictícios (Seção 8.3): 8 de 10 pares do mesmo animal sugeridos e 6 de 10 pares diferentes sugeridos por engano. \\ \hline
70\% dos participantes aprovaram a exigência de comprovação antes da devolução do animal, para prevenir golpes. & RF18 & Solicitação de devolução com relato e fotos, análise pelo autor da publicação, aprovação ou recusa, liberação do endereço exato após a aprovação e notificação do solicitante. & Fluxo principal verificado manualmente; sem erro observado. \\ \hline
70\% dos respondentes declararam que seus animais não têm microchip de identificação. & RF19 & Cadastro preventivo dos animais de estimação, carteirinha digital com código QR e contato do tutor, e conversão do cadastro em publicação de animal perdido. & Fluxo principal verificado manualmente; sem erro observado. \\ \hline
70\% dos respondentes afirmaram que usariam uma seção para adoção e lar temporário. & RF20 & Publicação para adoção e disponibilização de animal encontrado para adoção após sete dias de busca pelo tutor. & Fluxo principal verificado manualmente; sem erro observado. \\ \hline
\end{longtable}
}

\clearpage

% ============================================================
% APÊNDICE E: REGISTRO TÉCNICO DA ANÁLISE DE SEGURANÇA
% ============================================================
\section*{APÊNDICE E: REGISTRO TÉCNICO DA ANÁLISE DE SEGURANÇA}
\phantomsection
\label{apendice:analise-mobsf}
\addcontentsline{toc}{section}{APÊNDICE E: REGISTRO TÉCNICO DA ANÁLISE DE SEGURANÇA}

Este apêndice registra os dados necessários para identificar o pacote analisado e documenta as evidências produzidas pelo MobSF. O resumo SHA-256 identifica o arquivo específico submetido à análise; sua inclusão permite distinguir essa versão de outros pacotes com o mesmo nome.

\vspace{0.3cm}
\begin{table}[H]
\centering
\phantomsection\hypertarget{quadro:mobsf_pacote}{}\label{quadro:mobsf_pacote}
\small
\begin{tabular}{|>{\raggedright\arraybackslash}p{4.2cm}|>{\raggedright\arraybackslash}p{10.8cm}|}
\hline
\multicolumn{2}{|c|}{Quadro E1: Identificação do pacote submetido à análise estática} \\ \hline
Dado & Valor \\ \hline
Data da análise & 14 de outubro de 2026 \\ \hline
Ferramenta & Mobile Security Framework (MobSF), versão 4.5.4 \\ \hline
Tipo de análise & Estática; o aplicativo não foi executado \\ \hline
Arquivo analisado & \textit{WeFIND\_v2.0.apk} \\ \hline
Aplicativo e versão & WeFIND, versão 2.0 \\ \hline
Identificador do pacote & \textit{com.wefind.app} \\ \hline
Resumo SHA-256 & \parbox[t]{10.5cm}{\itshape\small f383cd85\allowbreak36142706\allowbreak d14d6e9d\allowbreak44bebdb2\allowbreak23ee7d78\allowbreak77bc9446\allowbreak a4fb21ff\allowbreak685b90f9} \\ \hline
\end{tabular}
\fonte{Relatório do MobSF (2026)}
\end{table}

\subsection*{Inspeção estática do código e das migrações}

Além da análise automatizada do APK, o código e as migrações do projeto foram inspecionados manualmente. Nessa inspeção, foi observado o uso de HTTPS na comunicação, com o tratamento de senhas delegado ao Supabase Auth e o uso de PKCE no fluxo OAuth. Também foram observadas políticas de segurança em nível de linha (RLS) nas tabelas que possuem regras definidas, além de validações de entrada. Os códigos temporários são gerados por fonte criptograficamente segura, apenas o resumo SHA-256 deles é armazenado e cada código aceita no máximo cinco tentativas. Por fim, a função de mensageria exige autenticação, e os dados sensíveis são removidos dos registros de log.

Essa inspeção estática não verificou de forma independente a infraestrutura hospedada, as configurações de backup e restauração, a expiração global de sessões ou a disponibilidade dos provedores utilizados. Também não foram realizados testes de invasão nem avaliação dinâmica do aplicativo.

\clearpage